\section{Privacy–Security Frontier: ampliar con ingeniería el espacio de opciones}

Palantir comenzó en el ámbito antiterrorista, por lo que desde su origen enfrenta controversias sobre la privacidad y las libertades civiles. Sankar describe la división de tareas entre la política y la ingeniería de la siguiente manera: la política decide en qué punto del privacy/security trade-off está dispuesta a situarse la sociedad, mientras que la ingeniería debe empujar la efficient frontier hacia afuera, de modo que con la misma protección de la privacidad se obtenga más seguridad, o que con la misma capacidad de seguridad se reduzca la intrusión en los datos. Este marco es útil, pero «empujar hacia afuera» debe sustentarse en mecanismos medibles y no convertirse en retórica para ampliar la vigilancia.

\subsection{El significado de la efficient frontier}

La \term{Efficient frontier} es la frontera en la que ya no es posible mejorar un objetivo sin sacrificar otro. Sea $P$ la protección de la privacidad y $S$ la capacidad de seguridad; la mejora de ingeniería busca que el conjunto factible se amplíe de $\mathcal{F}_0$ a $\mathcal{F}_1$, de modo que algunos puntos logren un $S$ mayor con el mismo $P$. Aquí cada símbolo requiere indicadores reales, como la minimización de datos, la tasa de falsos positivos, el tiempo de atención de casos y los accesos no autorizados, y no una abstracta «más seguridad».

\begin{figure}[H]
\centering
\includegraphics[width=0.94\textwidth]{images/12-privacy-security-frontier.png}
\caption{Las mejoras de ingeniería pueden empujar la privacy–security capability frontier, pero la sociedad aún debe elegir un punto aceptable.}
\figsource{Redibujado a partir de los subtítulos 39:30--42:00, `images/12-privacy-security-frontier.png`.}
\end{figure}

\begin{knowledgebox}{Lectura de la figura: qué mecanismos pueden empujar realmente la frontier}
El fine-grained access control reduce el número de personas ajenas que ven los datos; la purpose limitation restringe su uso; el immutable audit permite exigir responsabilidades por los abusos; la privacy-preserving linkage realiza el cruce de registros sin copiar todos los datos originales; la retention policy elimina automáticamente los datos vencidos. Cada mecanismo debe someterse a red-team, a pruebas de uso indebido y a supervisión independiente.
\end{knowledgebox}

\begin{warningbox}{La capability no sustituye el debido proceso}
Aunque el sistema sea técnicamente capaz de seleccionar a las personas con mayor precisión, sigue siendo necesario responder quién lo autoriza, quién puede apelar, cómo se corrigen los errores, cuánto tiempo se conservan los datos y qué pueden ver los supervisores. Mejorar la frontier amplía el espacio de opciones institucionales, pero no elige por sí sola una forma de uso legítima o justa.
\end{warningbox}

\subsection{Resumen de la sección}

Privacy y security no son un único control deslizante que solo pueda disputarse con consignas. El access control, la minimización de datos, la auditoría y la computación que preserva la privacidad pueden reducir parte del conflicto, pero los objetivos finales y la autorización siguen correspondiendo a la política y al derecho. La responsabilidad del equipo técnico es lograr que el trade-off sea medible, acotable y revisable.

